% Geometry and arithmetic of the weight-tripling correspondence curve C_3.

\begin{theorem}[三倍重量对应曲线的正规化与奇点 \texorpdfstring{$\delta$}{delta}-总亏格]\label{thm:xi-terminal-zm-elliptic-weight-tripling-c3-normalization-delta}
\leanverified{paper\_xi\_terminal\_zm\_elliptic\_weight\_tripling\_c3\_normalization\_delta}
沿用结论 \ref{con:xi-terminal-zm-elliptic-weight-tripling-modular-equation} 的记号。
令 \(M_3(y,y_3)\in\ZZ[y,y_3]\) 为满足 \(M_3(y(P),y([3]P))=0\) 的原始不可约多项式，且 \((\deg_y,\deg_{y_3})=(36,4)\)。
令
\[
C_3\subset \PP^1_y\times \PP^1_{y_3}
\]
为方程 \(M_3(y,y_3)=0\) 的射影闭包（取既约结构）。
则：
\begin{enumerate}
  \item \(C_3\) 的正规化同构于椭圆曲线 \(\mathcal E\)；等价地其函数域满足
  \[
  \QQ(C_3)\cong \QQ(\mathcal E).
  \]
  \item \(C_3\) 的算术亏格为
  \[
  p_a(C_3)=(36-1)(4-1)=105.
  \]
  令 \(S:=\mathrm{Sing}(C_3\otimes\overline{\QQ})\) 为几何奇点集，则
  \[
  \sum_{P\in S}\delta_P=p_a(C_3)-g(\mathcal E)=105-1=104.
  \]
\end{enumerate}
\end{theorem}
\begin{proof}
由不可约性与 \(\deg_{y_3}M_3=4\)，有
\(
[\QQ(y,y_3):\QQ(y)]=4
\)。
另一方面，\(y:\mathcal E\to\PP^1_y\) 为次数 \(4\) 的覆叠（结论 \ref{con:xi-terminal-zm-elliptic-y-hurwitz-passport}），故
\(
[\QQ(\mathcal E):\QQ(y)]=4
\)。
映射
\[
\mathcal E\dashrightarrow \PP^1_y\times\PP^1_{y_3},\qquad P\longmapsto (y(P),y([3]P))
\]
的像落在 \(C_3\) 上，并诱导函数域嵌入 \(\QQ(y,y_3)\hookrightarrow \QQ(\mathcal E)\)。
两侧作为 \(\QQ(y)\)-扩张的次数同为 \(4\)，故该嵌入为同构，从而 \(C_3\) 的正规化同构于 \(\mathcal E\)。
\par
对 \(\PP^1\times\PP^1\) 上双次数为 \((a,b)\) 的不可约曲线，其算术亏格由伴随公式给出 \(p_a=(a-1)(b-1)\)，取 \((a,b)=(36,4)\) 得 \(p_a(C_3)=105\)。
由正规化属数为 \(g(\mathcal E)=1\)，对平面曲线奇点的 \(\delta\)-不变量满足
\(
\sum_{P\in S}\delta_P=p_a(C_3)-g(\mathcal E)
\)，即得总和为 \(104\)。证毕。
\end{proof}

\begin{proposition}[判别式的场判别式与指数平方分解：Lee--Yang 分歧与非正规缺陷分离]\label{prop:xi-terminal-zm-elliptic-weight-tripling-discriminant-index}
\leanverified{paper\_xi\_terminal\_zm\_elliptic\_weight\_tripling\_discriminant\_index}
沿用结论 \ref{con:xi-terminal-zm-elliptic-weight-tripling-modular-equation} 的记号，并令 \(A:=\QQ[y]\)。
考虑次数 \(4\) 的 \(A\)-代数
\[
\mathcal O:=A[y_3]/(M_3),
\]
并令 \(K:=\mathrm{Frac}(\mathcal O)\cong \QQ(\mathcal E)\)。
记 \(\widetilde{\mathcal O}\) 为 \(\mathcal O\) 在 \(K\) 中的整闭包（等价地为 \(C_3\) 的正规化在 \(\mathbb A^1_y\) 上的仿射坐标环）。
则判别式满足标准恒等式
\[
\mathrm{Disc}_A(\mathcal O)=\mathrm{Disc}_A(\widetilde{\mathcal O})\cdot \mathrm{Index}(\widetilde{\mathcal O}:\mathcal O)^2,
\]
其中 \(\mathrm{Index}(\widetilde{\mathcal O}:\mathcal O)\subset A\) 为指数理想（忽略 \(A^\times=\QQ^\times\) 的单位因子）。
在本情形中可取
\[
\mathrm{Disc}_A(\widetilde{\mathcal O})=-y(y-1)P_{\mathrm{LY}}(y),
\qquad
\mathrm{Index}(\widetilde{\mathcal O}:\mathcal O)=(A_{32}(y)B_{66}(y)),
\]
并且与结论 \ref{con:xi-terminal-zm-elliptic-weight-tripling-disc-squareclass} 的显式分解一致：
\[
\mathrm{Disc}_{y_3}\bigl(M_3(y,y_3)\bigr)=-y(y-1)P_{\mathrm{LY}}(y)\cdot A_{32}(y)^2B_{66}(y)^2.
\]
从而 \(-y(y-1)P_{\mathrm{LY}}(y)\) 控制扩张 \(K/\QQ(y)\) 的场判别式（分歧）部分，而平方因子 \(A_{32}^2B_{66}^2\) 完全来自 \(\mathcal O\) 作为 \(K\) 的非极大序所产生的指数缺陷。
\end{proposition}
\begin{proof}
对 \(A\) 的适当局部化（去除有限集合的坏素）后，\(A\) 为 Dedekind 域且 \(\widetilde{\mathcal O}\) 为其上极大序。
判别式的基变换公式给出标准恒等式
\(
\mathrm{Disc}(\mathcal O)=\mathrm{Disc}(\widetilde{\mathcal O})\cdot[\widetilde{\mathcal O}:\mathcal O]^2
\)（函数域与数域同理）。
此处 \(\mathrm{Disc}_{y_3}(M_3)\) 为 \(\mathcal O\) 的多项式判别式生成元（差一个单位）。
\par
另一方面，\(\widetilde{\mathcal O}\) 的判别式理想由扩张 \(K/\QQ(y)\) 的分歧决定，而该分歧集合与四次覆盖 \(y:\mathcal E\to\PP^1_y\) 的分歧值一致；
由结论 \ref{con:xi-terminal-zm-elliptic-y-hurwitz-passport}，其分歧值恰由 \(-y(y-1)P_{\mathrm{LY}}(y)\) 锁定，故可取 \(\mathrm{Disc}_A(\widetilde{\mathcal O})=-y(y-1)P_{\mathrm{LY}}(y)\)（差一个单位）。
将其与结论 \ref{con:xi-terminal-zm-elliptic-weight-tripling-disc-squareclass} 的显式分解比较，即得指数理想可由 \(A_{32}B_{66}\) 生成（亦差一个单位）。证毕。
\end{proof}

\begin{theorem}[奇点投影刚性与近乎 \texorpdfstring{$\delta=1$}{delta=1} 的定量界]\label{thm:xi-terminal-zm-elliptic-weight-tripling-singularity-projection}
\leanverified{paper\_xi\_terminal\_zm\_elliptic\_weight\_tripling\_singularity\_projection}
沿用定理 \ref{thm:xi-terminal-zm-elliptic-weight-tripling-c3-normalization-delta} 与命题 \ref{prop:xi-terminal-zm-elliptic-weight-tripling-discriminant-index} 的记号。
记 \(S:=\mathrm{Sing}(C_3\otimes\overline{\QQ})\) 为几何奇点集，并令 \(\pi_1:C_3\to\PP^1_y\) 为第一投影。
则：
\begin{enumerate}
  \item 集合论支撑满足
  \[
  \pi_1(S)\subseteq Z(A_{32}B_{66})\subset \PP^1_y.
  \]
  并且 \(Z(A_{32}B_{66})\) 与 Lee--Yang 分歧集合 \(Z(y(y-1)P_{\mathrm{LY}})\) 互不相交。
  \item \(A_{32}\) 与 \(B_{66}\) 在 \(\overline{\QQ}[y]\) 中均平方自由，且 \(\gcd(A_{32},B_{66})=1\)。
  因此
  \[
  \#Z(A_{32}B_{66})=32+66=98.
  \]
  \item 由 \(\sum_{P\in S}\delta_P=104\) 且 \(\delta_P\ge 1\)（任意曲线奇点），得到
  \[
  98\le \#S\le 104,
  \qquad
  \sum_{P\in S}(\delta_P-1)=104-\#S\le 6.
  \]
  特别地，满足 \(\delta_P\ge 2\) 的奇点至多 \(6\) 个，其余奇点皆为 \(\delta=1\) 型。
\end{enumerate}
\end{theorem}
\begin{proof}
\textbf{(1)}
由命题 \ref{prop:xi-terminal-zm-elliptic-weight-tripling-discriminant-index}，\(\mathcal O\) 的非正规 locus 在底环 \(A=\QQ[y]\) 上的支撑由指数理想 \((A_{32}B_{66})\) 控制。
对一维既约曲线，仿射坐标环在点处正规当且仅当该点为非奇点；因此奇点在 \(y\)-线上之像包含于 \(Z(A_{32}B_{66})\)。
由脚本 \texttt{scripts/exp\_xi\_weight\_tripling\_index\_squarefree\_audit.py} 的可复现核验，\(\gcd(A_{32}B_{66},y(y-1)P_{\mathrm{LY}})=1\)，故非正规支撑与分歧支撑互不相交。
\par
\textbf{(2)}
同一核验脚本给出 \(\gcd(A_{32},A_{32}')=\gcd(B_{66},B_{66}')=\gcd(A_{32},B_{66})=1\)，从而 \(A_{32},B_{66}\) 皆平方自由且互素。
于是零点不重数且无交叉，故 \(\#Z(A_{32}B_{66})=\deg A_{32}+\deg B_{66}=98\)。
\par
\textbf{(3)}
由 (1)，每个 \(y\in Z(A_{32}B_{66})\) 的纤维 \(\pi_1^{-1}(y)\) 上至少出现一个非正规点，从而 \(\#S\ge \#Z(A_{32}B_{66})\)。
另一方面由 \(\delta_P\ge 1\) 得 \(\#S\le \sum_{P\in S}\delta_P=104\)。
最后
\(
\sum_{P\in S}(\delta_P-1)=\sum_{P\in S}\delta_P-\#S
\)
给出上界 \(104-\#S\le 6\)，并推出 \(\delta_P\ge 2\) 的奇点至多 \(6\) 个。证毕。
\end{proof}

\begin{corollary}[乘三下 Lee--Yang 分歧点的 \texorpdfstring{$E[3]$}{E[3]}-余类复制与精确分歧计数]\label{cor:xi-terminal-zm-elliptic-weight-tripling-ramification-count}
\leanverified{paper\_xi\_terminal\_zm\_elliptic\_weight\_tripling\_ramification\_count}
设 \(y:\mathcal E\to\PP^1_y\) 为次数 \(4\) 的重量映射，并令 \(y_3:=y\circ[3]\)。
记 \(R(y)\) 与 \(R(y_3)\) 为相应的分歧除子。
则在特征 \(0\) 下倍点同态 \([3]\) 为 \'etale，因而
\[
R(y_3)=[3]^*R(y).
\]
在结论 \ref{con:xi-terminal-zm-elliptic-y-hurwitz-passport} 的分歧刻画下，可写为
\[
R(y_3)=3\sum_{Q\in \mathcal E[3]}[Q]\ +\ \sum_{i=1}^{5}\ \sum_{Q\in \mathcal E[3]}[P_i+Q],
\]
其中 \(P_1,\dots,P_5\) 为 \(y\) 的五个有限分歧点。
从而分歧度具有精确计数：
\[
\deg R(y_3)=9\cdot(4-1)+45\cdot(2-1)=27+45=72,
\]
并与 \(\deg(y_3)=\deg(y)\deg([3])=4\cdot 9=36\) 的 Hurwitz 公式严格一致。
\end{corollary}
\begin{proof}
在特征 \(0\) 下，\([3]\) 处处非分歧，故对任意有理函数 \(f\) 有 \(R(f\circ[3])=[3]^*R(f)\)。
由结论 \ref{con:xi-terminal-zm-elliptic-y-hurwitz-passport}，\(\infty\) 为全分歧值且 \(e_O(y)=4\)，五个有限分歧点均为单分歧（指数 \(2\)）。
又 \([3]^{-1}(O)=\mathcal E[3]\)，且 \([3]^{-1}(P_i)=P_i+\mathcal E[3]\)。
由于 \'etale 拉回保持分歧指数，\(\mathcal E[3]\) 上各点对 \(y_3\) 的贡献为 \(4-1=3\)，而每个 \(P_i+\mathcal E[3]\) 上各点贡献为 \(2-1=1\)。
合计得到 \(\deg R(y_3)=9\cdot 3+5\cdot 9\cdot 1=72\)。
最后由 Riemann--Hurwitz（\(g(\mathcal E)=1\)）得 \(\deg R(y_3)=2\deg(y_3)=72\)，与计数一致。证毕。
\end{proof}

\endinput
